% 张三慧《学习辅导与习题解答》PDF S 独立选答，卷1第1章；未取得V1主教材独立交叉核验。
% 仅录“思考题选答”中实际明确列出的原编号；其余不推断，exam_ready=false。
\section*{第1章\quad 质点运动学：辅导书已列思考题独立解析}
\paragraph{思考题 1.7}
万有引力始终指向太阳焦点。按开普勒第二定律，单位质量角动量 \(h=r^2\dot\theta\) 守恒。在近日点 \(r_p\) 最小，\(v_p=h/r_p\) 最大；在远日点 \(r_a\) 最大，\(v_a=h/r_a\) 最小。相应 \(v_p/v_a=r_a/r_p>1\)。在其他轨迹位置需由速度切向分量及轨道能量共同确定，不能简单把加速度与速度同向画在全轨道。

\paragraph{思考题 1.8}
最高点速度仍为水平的 \(v_x=v_0\)，重力加速度竖直向下，与速度垂直，完全为法向加速度。用 \(a_n=v^2/\mathcal R\)，得 \(\boxed{\mathcal R=v_0^2/g}\)。这一结论只针对最高点与不计阻力的平抛/斜抛；其他位置需将 \(g\) 分解为切向和法向分量。

\paragraph{思考题 1.11}
取同一轨迹 \(x'(t')=x(-t')\)，则 \(\boxed{v'(t')=-v(-t')}\)、\(\boxed{a'(t')=a(-t')}\)。时间反演令速度反向而加速度保持同向，故倒放自由落体表现为上抛后减速，仍符合恒定向下重力加速度。在只受位置依赖的保守力时运动方程具有这种时间反演对称性；有摩擦等耗散作用时不成立。
